% Generated draft: requires mathematical and source review.
\documentclass[11pt]{article}
\usepackage[margin=25mm]{geometry}
\usepackage{fontspec}
\setmainfont{Latin Modern Roman}
\usepackage{amsmath,amssymb,mathtools,unicode-math}
\setmathfont{Latin Modern Math}
\usepackage{enumitem,needspace,xurl,hyperref}
\setlength{\parindent}{0pt}
\setlength{\parskip}{5pt}
\setlength{\emergencystretch}{3em}
\newcommand{\N}{\mathbb{N}}
\newcommand{\Z}{\mathbb{Z}}
\newcommand{\Q}{\mathbb{Q}}
\newcommand{\R}{\mathbb{R}}
\newcommand{\C}{\mathbb{C}}
\newcommand{\F}{\mathbb{F}}
\newcommand{\A}{\mathbb{A}}
\newcommand{\PP}{\mathbb{P}}
\newcommand{\E}{\mathbb{E}}
\newcommand{\eps}{\varepsilon}
\newcommand{\dd}{\,\mathrm{d}}
\begin{document}
\section*{UnsolvedMath ID: 30003158}
\textbf{Source problem code:} catalogue entry
\subsection{Definitions and mathematical statement}
Let \(K\) be a number field, let \(S\) be a finite, possibly empty, set of finite places of \(K\), and let
\[
\mathcal O_K[S^{-1}]=\{a\in K:v(a)\geq 0\text{ for every finite place }v\notin S\},\qquad B=\operatorname{Spec}(\mathcal O_K[S^{-1}]).
\]
Under the usual geometric-fiber interpretation, a smooth Fano threefold over \(B\) is a smooth projective morphism \(f:X\to B\) whose geometric fibers are connected three-dimensional varieties and whose relative anticanonical line bundle \(\omega_{X/B}^{-1}\) is relatively ample. For a geometric fiber \(X_{\bar b}\), write
\[
\rho(X_{\bar b})=\dim_{\Q}\bigl(\operatorname{NS}(X_{\bar b})\otimes_{\Z}\Q\bigr),
\]
where \(\operatorname{NS}\) denotes the group of line bundles modulo algebraic equivalence. In the formulation below, Picard rank one means \(\rho(X_{\bar b})=1\) for every geometric point \(\bar b\to B\). The supplied record does not explicitly specify whether Picard rank is geometric or arithmetic, or whether it is imposed on every fiber or only on the generic fiber; this convention therefore requires confirmation.

Two such families \(f:X\to B\) and \(f':X'\to B\) are \(B\)-isomorphic if there is an isomorphism \(\phi:X\to X'\) satisfying \(f'\circ\phi=f\). No marking or chosen polarization is included in this equivalence relation. Let \(\mathcal F_1(B)\) be the set of these \(B\)-isomorphism classes.

\textbf{General finiteness question.} Is it true that, for every number field \(K\) and every finite set \(S\) of finite places of \(K\),
\[
\#\mathcal F_1\bigl(\operatorname{Spec}(\mathcal O_K[S^{-1}])\bigr)<\infty?
\]

\textbf{Distinguished Gushel--Mukai case.} Does the same finiteness assertion hold for the subclass consisting of smooth Gushel--Mukai Fano threefolds over \(B\)? To identify the characteristic-zero geometric objects, a smooth Gushel--Mukai threefold over an algebraically closed field \(k\) of characteristic zero is a smooth three-dimensional intersection
\[
C\operatorname{Gr}(2,V_5)\cap\mathbf P(W)\cap Q
\subseteq\mathbf P(k\oplus\textstyle\bigwedge^2 V_5),
\]
where \(\dim_k V_5=5\), \(C\operatorname{Gr}(2,V_5)\) is the projective cone over the Pl\"ucker-embedded Grassmannian, \(W\subseteq k\oplus\bigwedge^2 V_5\) has dimension eight, and \(Q\) is a quadric hypersurface in \(\mathbf P(W)\); the intersection avoids the cone vertex. This description identifies the geometric type but does not settle the intended definition over \(B\). The record does not specify whether Gushel--Mukai type is required only on the geometric generic fiber or on all geometric fibers, nor does it supply an integral definition at residue characteristic two. Consequently, the restricted finiteness question retains this unresolved qualification. No restriction on residue characteristics is stated in the source.

The original statement singles out the Gushel--Mukai case, whereas the cleaned statement also asks the general Picard-rank-one question. Finiteness for the general class would imply finiteness for its Gushel--Mukai subclass; a converse is not asserted here.
\subsection{Short English statement}
For every number field and every finite set of allowed finite places, ask whether only finitely many smooth Fano threefold families of Picard rank one exist over the corresponding ring of \(S\)-integers, up to isomorphism over that ring. The source particularly asks this for Gushel--Mukai threefolds. The Picard-rank convention and the integral meaning of Gushel--Mukai type require clarification. This is a statement-only draft requiring human mathematical review.
\subsection{Sources}
\begin{itemize}
\item Local UnsolvedMath registry: supplied complete record with permanent ID \(30003158\), exact source problem code \texttt{catalogue entry}, title ``Arithmetic Finiteness of Gushel--Mukai Threefolds,'' in the Oberwolfach Reports Open Problems collection. No individual proposer is supplied; the proposed year is recorded as 2016.
\item Supplied report citation: \emph{Classical Algebraic Geometry (2017)}, Oberwolfach Reports, DOI \texttt{10.4171/owr/2016/29}, pp.\ 1635--1682. \url{https://doi.org/10.4171/owr/2016/29}. The background attributes the problem to workshop open-problem discussions. The citation's 2017 designation and the DOI's 2016 component are retained as supplied.
\item Supplied literature source: ``Oberwolfach Report 29.'' \url{https://doi.org/10.4171/owr/2016/29}.
\item Supplied literature source: ``Gushel--Mukai moduli with level structures.'' \url{https://arxiv.org/abs/1511.09291}. The title and link are reproduced as supplied; their correspondence has not been independently verified.
\item Literature assessment dated \texttt{2026-08-21}: the supplied curation labels the problem open and states that finiteness of \(B\)-isomorphism classes of smooth Gushel--Mukai Fano threefolds remains unknown in the report, with only conditional Lang--Vojta-type finiteness available. This is supplied dated curation, not an independently verified literature assessment; its time-sensitive status and cited claims require independent verification.
\item The registry labels the statement \texttt{corrected\_verified} and says it was checked against the official report. These are supplied verification claims, not verification performed for this draft.
\end{itemize}
\end{document}
